\subsection{Discussão com a literatura}
\label{sec:discussao-literatura}
Os estudos relacionados ajudam a interpretar mecanismos, mas não fornecem uma escala comum para comparar milissegundos ou percentuais. Linguagem, topologia, carga, significado de resposta e implantação diferem. Por isso, confrontam-se direções e limites, sem somar as medições de aplicações distintas como se fossem repetições deste experimento.

\citeonline{cabane2024} compararam uma aplicação monolítica com sua alternativa orientada a eventos. O manuscrito disponibilizado pelos autores relata maior uso de memória na alternativa distribuída e discute o custo de executar múltiplas subaplicações. A direção é compatível com a maior memória do caminho assíncrono aqui observada, que inclui broker e consumidor adicionais. Isso não identifica a causa única nem transforma a topologia em medida de custo financeiro. O estudo também mostra por que a definição de throughput importa: respostas entre módulos podem participar da contagem utilizada em uma aplicação distribuída. Neste trabalho, vazão é restrita a registros únicos concluídos no banco durante a janela de medição, e não à soma de mensagens ou respostas internas.

A comparação do tempo de resposta exige outra distinção. O sucesso da API síncrona deste protótipo sucede a gravação; o aceite assíncrono sucede apenas a publicação. O menor p95 HTTP assíncrono, acompanhado de maior tempo até a gravação, não contradiz estudos nos quais ``resposta'' corresponde ao término de uma função mais abrangente. Também não permite afirmar que EDA seja sempre mais rápida. As cargas e as unidades de trabalho precisam ser equivalentes antes de comparar a direção de uma diferença.

\citeonline{kazanavicius2023} examinaram diferentes mecanismos de comunicação em um contexto de decomposição de monólitos. A seleção depende das necessidades da aplicação e das características dos mecanismos, não apenas de uma ordenação geral de desempenho. O resultado deste TCC reforça essa leitura contextual: se o requisito é saber imediatamente que o registro foi gravado, um aceite antecipado não atende à mesma finalidade; se é continuar recebendo durante a parada do consumidor, a retenção intermediária observada é relevante. O experimento não reproduz uma migração de monólito nem avalia as demais alternativas desse estudo.

\citeonline{blinowski2022} mostraram a importância do ambiente e do dimensionamento ao comparar monólitos e microsserviços em implementações e implantações distintas. Acrescentar distribuição ou instâncias não assegura ganho proporcional. Aqui, o menor uso de recursos do caminho síncrono e a ausência de vantagem de vazão nas taxas oferecidas devem ser interpretados nessa mesma cautela. Não houve busca da capacidade máxima, nem variação do número de brokers ou consumidores; portanto, atingir a mesma taxa não demonstra capacidade equivalente, e maior CPU não demonstra menor escalabilidade em outras configurações.

Os destinos durante falha dialogam com a discussão de \citeonline{helland2007} sobre trabalho fora de uma transação global e com os desafios de gestão de dados caracterizados por \citeonline{laigner2021}. Confirmar publicação, persistir no banco e confirmar a posição de consumo são passos diferentes. Idempotência por identificador protege o efeito persistido contra certas reentregas, mas não cria atomicidade automática entre Kafka e PostgreSQL. Os eventos encaminhados à DLQ continuaram identificáveis para inspeção; contudo, sem reprocessamento, não atingiram o objetivo de gravação da auditoria. A distinção entre retenção e persistência é uma consequência operacional concreta dessa separação de responsabilidades.

\citeonline{niedermaier2019} oferecem evidência qualitativa sobre desafios de observabilidade de sistemas distribuídos, não um benchmark para comparar nossas latências. A coleta temporal incompleta ilustra a importância dessa preocupação: a linha encontrada no banco permitiu verificar conteúdo e destino, mas não recuperar seu tempo de conclusão. Logs correlacionados, inventários e conciliação tornaram a limitação detectável; eles não a eliminaram. Da mesma forma, a análise exploratória de recuperação mostrou que a definição de pendência pode mudar a interpretação de um indicador composto. O critério original foi preservado para manter explícita essa diferença.

Em conjunto, a literatura e o experimento sustentam avaliar a função realmente confirmada, os componentes executados e os destinos de falha. O trabalho acrescenta uma comparação local com função de registro comum e marcos separados de aceite e conclusão; não demonstra superioridade universal, tolerância a desastre ou garantia de processamento exatamente uma vez.
